\section{考察}
\label{sec_discussion}
\subsection{ラベルの定義と精度の上限}
本研究で最も大きな影響を持ったのは，ラベルの定義であった．
必要な知識の4段階は，目視では上から順に判定でき，判断の手順として妥当であったが，自動判定のマクロ F1 は0.48を超えなかった．
2値に改めて条件を厳密にしたことで，教師の画像外の知識は F1 0.32から0.67に上がった．
これは主に，指示に依存する事物の種類を選ばせ，普通名詞を特定の事物とする誤りを抑えたことによる．
また，付け直しによる再現性の測定は，同じ判断者でも6区分の約14\%，画像外の知識の約7\%の判断が揺れ，F1 ではそれぞれ0.80と0.85にとどまることを示した．
自動判定の目標値を決める前に，この上限を測っておくことが重要である．
一方で，この上限は同一の判断者の再現性であり，複数の判断者の間の一致度はさらに低い可能性がある．

\subsection{教師の偏りと少量の目視ラベル}
VLM の教師は，セリフをツッコミとし，状況の説明をほとんど使わず，一般的な語を特定の事物とするといった系統的な偏りを持っていた．
蒸留した生徒はこの偏りを受け継ぐため，教師ラベルだけでは教師の水準を超えにくい．
指示の分解や例示による偏りの修正は，かえって別の偏りを生んだ．
これに対し，931件の目視ラベルによる微調整は，6区分を0.55から0.71に上げた．
大量の教師ラベルで画像外の知識に必要な表現を学び，少量の目視ラベルで区分の境界を学ぶ二段階の学習は，教師の偏りを直接直すより効率的であった．
ただし，目視ラベルを増やす効果は緩やかであり，6区分で0.8に近づけるには，目視ラベルを大きく増やすか，揺れの大きいツッコミと心の声の定義を見直す必要がある．

\subsection{画像外の知識の限界}
画像外の知識の見逃しの多くは，有名な台詞や定型の引用に気づかないことによる．
字面が普通の日本語である引用は，元ネタを知らなければ判定できず，2〜8B のモデルの持つ知識では限界がある．
32B の教師も元ネタを書き出させると過剰に「はい」とし，教師ラベルの改善には結びつかなかった．
より大きなモデルを目視ラベルで微調整して教師とすることや，お題画像ごとに写っている作品を一度だけ同定して入力に加えることが，次の候補である．
後者はお題ごとに1回の処理であり，\ref{sec_speed}章のプレフィックスの共有とも両立する．

\subsection{本研究の限界}
本研究の目視ラベルは，全て AI エージェントである Claude Code が付与したものであり，人の判断者によるラベルとの一致は測っていない．
大喜利の解釈には文化的な知識や感覚が関わるため，人の判断者による検証が必要である．
また，検証データは620件であり，最終評価用データの F1 の95\%信頼区間は約 $ \pm 0.08 $ と広い．
本稿で比較した多くの改善は，この幅より小さい．
さらに，本稿の方法はまだ全件には適用しておらず，全件のラベルの分布や，下流の生成・判定での有用性は検証していない．
